\chapter{进程调度}

\section{调度基本概念}
\concept{调度队列}
\begin{points}
  \item 作业队列：外存上等待进入系统的作业集合。
  \item 就绪队列：内存中已就绪、等待 CPU 的进程集合。
  \item 设备队列：等待某个 I/O 设备的进程集合。
  \item 进程在整个生命周期中会在这些队列之间迁移。
\end{points}
\plain{这三个队列对应三种典型等待对象：等进入系统、等 CPU、等设备。只要先看“在等谁”，队列就不容易混。}
\exam{选择题很喜欢把“外存后备作业”“内存就绪进程”“等待 I/O 的阻塞进程”三类对象打乱让你配对。}


\concept{三级调度}
\begin{points}
  \item 高级调度又称作业调度/长程调度，决定哪些作业从外存进入内存，影响多道程序度。
  \item 中级调度又称交换调度，决定哪些进程换入、换出内存，用于缓解内存紧张。
  \item 低级调度又称进程调度/短程调度，决定就绪队列中哪个进程获得 CPU，运行频率最高。
  \item 记忆：高级管“进不进系统”，中级管“在不在内存”，低级管“谁上 CPU”。
\end{points}
\plain{三级调度最容易记混的是“管到哪一层”：高级决定进系统，中级决定留不留内存，低级决定此刻谁拿 CPU。}
\exam{常考三个调度层次的别名和作用范围，尤其中级调度 = 交换调度，别漏。}


\concept{调度要解决的三个问题}
\begin{points}
  \item \term{WHAT}：按什么原则选择下一个要执行的进程，即调度算法本身。
  \item \term{WHEN}：何时进行调度，即调度时机，和抢占式/非抢占式密切相关。
  \item \term{HOW}：怎样把被选中的进程真正放上 CPU，即分派和上下文切换过程。
\end{points}
\plain{这页 PPT 的主线很适合简答题：先说选谁，再说什么时候选，最后说怎么切过去。}
\exam{常考把调度问题拆成 WHAT/WHEN/HOW，或据此解释算法、时机和切换过程的区别。}


\concept{CPU 密集型与 I/O 密集型}
\begin{points}
  \item CPU 密集型作业计算时间多，CPU burst 长，I/O 次数相对少。
  \item I/O 密集型作业 I/O 操作多，CPU burst 短。
  \item 作业调度应合理搭配两类作业：CPU 密集型过多会导致 I/O 设备空闲；I/O 密集型过多会导致就绪队列空、CPU 空闲。
\end{points}
\plain{调度不仅看“谁先跑”，还要看作业类型怎么搭配，才能让 CPU 和 I/O 设备都别闲着。}
\exam{常考 CPU 密集型和 I/O 密集型区别，以及为什么作业调度要混合搭配这两类作业。}


\concept{调度程序与分派程序}
\begin{points}
  \item 调度程序负责从就绪队列中选择下一个运行进程。
  \item 分派程序负责完成上下文切换、切换到用户态、跳转到被选进程的正确位置。
  \item 分派延迟是停止一个进程并启动另一个进程所需时间，应尽可能短。
\end{points}
\plain{调度器负责“选人”，分派器负责“把 CPU 真交过去”。}
\exam{常考 scheduler 和 dispatcher 的分工，以及什么叫分派延迟。}


\concept{上下文切换}
\begin{points}
  \item 上下文切换是保存当前进程状态并恢复另一个进程状态的过程。
  \item 上下文信息通常保存在 PCB 中。
  \item 上下文切换时间属于纯系统开销，不直接创造用户计算结果。
  \item 因此时间片过短、抢占过于频繁时，调度开销会明显上升。
\end{points}
\plain{切换本身不干活，只是在“换人上机”，所以切得太勤会把 CPU 时间浪费在管理上。}
\exam{常考上下文切换保存到哪里、为什么它属于系统开销、与时间片大小有什么关系。}


\section{调度目标与评价指标}
\concept{常见评价指标}
\begin{points}
  \item CPU 利用率：CPU 忙碌时间比例，越高越好。
  \item 吞吐量：单位时间完成的进程数量，越高越好。
  \item 周转时间：进程从提交到完成的总时间，越短越好。
  \item 等待时间：进程在就绪队列中等待 CPU 的总时间，越短越好。
  \item 响应时间：从提交请求到首次响应的时间，分时系统尤其关注。
  \item 公平性：不同用户或进程应获得合理的 CPU 份额，避免长期饥饿；不同系统对“公平”的定义并不完全相同。
\end{points}
\plain{这些指标不是并列硬背，真正的源头通常只有调度时间轴；时间轴一错，后面周转、等待、响应会全错。}
\exam{计算题通常先让你画甘特图，再顺着求完成时间和各指标；只写结果不写过程很容易丢步骤分。}


\concept{周转时间公式}
\begin{points}
  \item 完成时间 $C_i$，到达时间 $A_i$，服务时间 $S_i$。
  \item 周转时间 $T_i=C_i-A_i$。
  \item 带权周转时间 $W_i=\dfrac{T_i}{S_i}$。
  \item 平均周转时间和平均带权周转时间就是对所有进程取平均。
\end{points}
\plain{公式本身不难，难点是先把 $C_i$ 算对；周转和带权周转基本都是从完成时间往后顺推。}
\exam{常见陷阱是把服务时间当周转时间分母之外的别的量，或把到达时间不是 0 的进程直接套错公式。}


\concept{等待时间与响应时间}
\begin{points}
  \item 等待时间是进程在就绪队列中等待 CPU 的总时间。
  \item 对非抢占式算法，等待时间常可由“周转时间减服务时间”得到。
  \item 响应时间是从提交请求到第一次得到响应的时间，交互系统特别关心。
  \item 时间片轮转等分时算法常优先优化响应时间，而不一定使平均周转时间最小。
\end{points}
\plain{批处理更关心“多久做完”，交互系统更关心“多久先有反应”。}
\exam{常考等待时间、周转时间、响应时间的区别，不要把“第一次响应”和“最终完成”混掉。}


\concept{调度题作答顺序}
\begin{points}
  \item 先按题目给出的到达时间、服务时间、优先级、时间片等规则画出调度顺序或甘特图。
  \item 再列每个进程的完成时间。
  \item 然后依次计算周转时间、等待时间、带权周转时间等指标。
  \item 只写最终数字而不写过程，考试里很容易丢步骤分。
\end{points}
\plain{老师录音特别强调这类题要把过程写出来，因为后面所有指标都依赖前面的调度顺序。}
\exam{计算题中通常既看结果也看过程，至少要把调度序列和每个指标的来源写清楚。}


\section{调度方式}
\concept{抢占式与非抢占式}
\begin{points}
  \item 非抢占式调度：进程一旦获得 CPU，除非主动放弃、阻塞或结束，否则不会被强行剥夺。
  \item 抢占式调度：操作系统可因时间片到、更高优先级进程到达等原因剥夺 CPU。
  \item 分时系统通常需要抢占式调度；批处理系统可以使用非抢占式调度。
\end{points}
\plain{抢占就是正在运行的进程也可能被更高优先级/时间片到期打断。}
\exam{常考抢占式与非抢占式适用场景。}


\concept{调度发生的时机}
\begin{points}
  \item 非抢占式调度通常在进程终止，或进程由运行态转入阻塞态时发生。
  \item 抢占式调度除上述时机外，还可在时间片到、运行态转就绪态、高优先级进程就绪、新进程进入就绪队列等情况下发生。
  \item 记忆关键：抢占式比非抢占式多了“即使当前进程还没结束，也可能因为外界条件变化被换下”。
\end{points}
\plain{调度时机题本质是在问：当前正在跑的进程，什么时候必须交棒，什么时候可能被强行收回 CPU。}
\exam{常考非抢占和抢占条件对比，尤其时间片到和高优先级进程到达是否触发调度。}


\section{常用调度算法}
\concept{先来先服务 FCFS}
\begin{points}
  \item 按进程到达就绪队列的先后顺序分配 CPU。
  \item 优点：简单、公平、不会饥饿。
  \item 缺点：短作业可能排在长作业后面，平均等待时间可能很长，存在“护航效应”。
  \item 适合批处理场景，不适合交互式短任务较多的场景。
\end{points}
\plain{谁先来谁先上，最朴素但容易让短作业排在长作业后面。}
\exam{常考护航效应和平均等待时间计算。}


\concept{短作业优先 SJF 与最短剩余时间优先 SRTF}
\begin{points}
  \item SJF 选择预计运行时间最短的作业先执行，是非抢占式算法。
  \item SRTF 是抢占式版本，总是选择剩余运行时间最短的进程。
  \item 优点：理论上可使平均等待时间最短。
  \item 缺点：需要预测运行时间；长作业可能长期等待，产生饥饿。
\end{points}
\plain{优先做短的，平均等待时间好看，但长作业可能一直等。}
\exam{常考非抢占 SJF 和抢占式 SRTF 的甘特图区别。}


\concept{SJF 为什么平均等待时间最短}
\begin{points}
  \item 在已知运行时间且都同时到达的理想条件下，SJF 可使平均等待时间达到最小。
  \item 直观原因是先完成短任务，能减少后续所有任务累计等待。
  \item 但现实系统通常无法准确预知服务时间，因此常需基于历史行为估计。
\end{points}
\plain{短任务先做，后面排队的人整体少等一些，所以平均值最好看。}
\exam{常考“哪种算法平均等待时间理论最优”，答案通常是 SJF/SRTF。}


\concept{优先级调度}
\begin{points}
  \item 为每个进程赋予优先级，优先级高者先运行。
  \item 可分为抢占式和非抢占式。
  \item 静态优先级创建后不变，动态优先级可根据等待时间、资源使用情况改变。
  \item 缺点是低优先级进程可能饥饿；常用老化 aging 提高等待过久进程的优先级。
\end{points}
\plain{给进程分等级，重要的先跑；问题是低优先级可能饥饿。}
\exam{常考静态/动态优先级、抢占/非抢占、老化技术防饥饿。}


\concept{时间片轮转 RR}
\begin{points}
  \item 每个就绪进程轮流获得一个时间片，时间片用完仍未完成则回到就绪队列尾部。
  \item 适合分时系统，强调响应时间和公平性。
  \item 时间片太大，退化为 FCFS；时间片太小，上下文切换开销过大。
  \item 做题时严格画时间轴，注意新到达进程插入队列的位置和时间片剩余规则。
\end{points}
\plain{每人轮流跑一小段，适合分时系统；时间片太大像 FCFS，太小切换开销大。}
\exam{常考时间片轮转甘特图，以及时间片大小对响应时间和开销的影响。}


\concept{高响应比优先 HRRN}
\begin{points}
  \item 设等待时间为 $W$、服务时间为 $S$，则响应比公式为 $R=\dfrac{W+S}{S}=1+\dfrac{W}{S}$。
  \item 等待时间越长，响应比越大；服务时间越短，响应比也越大。
  \item 该算法兼顾短作业优先和防止长作业饥饿。
  \item 通常是非抢占式调度，用在作业调度中较常见。
\end{points}
\plain{它像给短作业加分，同时又让等太久的长作业慢慢把分数追上来。}
\exam{常考响应比公式计算，以及 HRRN 为什么能缓解 SJF 的饥饿问题。}


\concept{多级队列调度}
\begin{points}
  \item 将就绪队列划分为多个队列，不同队列服务不同类型进程，如系统进程、交互进程、批处理进程。
  \item 队列之间有固定优先级或时间比例分配，队列内部可使用不同调度算法。
  \item 固定多级队列的问题是低优先级队列可能长期得不到服务。
\end{points}
\plain{先把不同类型任务分到不同窗口，再决定窗口之间谁优先。}
\exam{常考多级队列和多级反馈队列的区别：前者队列固定，后者进程可在队列间移动。}


\concept{多级反馈队列 MLFQ}
\begin{points}
  \item 多级反馈队列允许进程在不同优先级队列之间移动。
  \item 新进程通常先进入高优先级队列，时间片较短，若用完时间片还未完成则降级。
  \item 等待过久的低优先级进程可被提升，避免饥饿。
  \item 直觉：先假设新来的任务可能是短交互任务，给它快速响应；若它表现得像长计算任务，就逐步降到低优先级。
\end{points}
\plain{这是老师点名的重点算法：先快照顾短交互任务，再把长任务慢慢往低队列放。}
\exam{常考调度大题：写清抢占规则、降级规则、时间片大小、低优先级进程是否会老化提升。}


\section{实时调度、多处理器调度与线程调度}
\concept{实时调度基本概念}
\begin{points}
  \item 实时系统中，正确结果如果迟到，也可能等同于错误结果。
  \item 周期性任务常用处理时间 $t_i$、截止期限 $d_i$、周期 $p_i$ 描述，通常 $0\le t_i\le d_i\le p_i$。
  \item 必要的可调度条件之一：$\sum_{i=1}^{m}\dfrac{t_i}{p_i}\le 1$，表示 CPU 总利用需求不能超过 100\%。
\end{points}
\plain{这里的利用率条件是在做“算术上装不装得下”的第一层判断：若总需求都超过 100\%，后面再好的算法也救不了。}
\exam{常考给出若干周期任务的 $t_i,p_i$，先判断总利用率是否至少满足可调度必要条件。}


\concept{硬实时、软实时与两类延迟}
\begin{points}
  \item 硬实时系统存在必须满足的时间限制，错过截止期限通常不可接受。
  \item 软实时系统偶尔超时可以容忍，但仍希望尽量及时。
  \item \term{中断延迟}是从 CPU 收到中断到中断处理程序开始的时间。
  \item \term{分派延迟}是调度程序从停止一个进程到切换到另一个进程的时间。
  \item 实时系统要求这两类延迟都尽量短，且上界可预测。
\end{points}
\plain{实时系统不只看算法名字，往往先考“什么叫中断延迟、什么叫分派延迟，它们为什么要可预测”。}
\exam{常考硬实时/软实时区别，以及中断延迟和分派延迟分别指什么。}


\concept{速率单调 RMS 与最早截止期限 EDF}
\begin{points}
  \item RMS 是静态优先级算法，周期越短，优先级越高。
  \item EDF 是动态优先级算法，截止期限越近，优先级越高。
  \item EDF 在单处理器、理想条件下可达到更高利用率，但实现和开销更复杂。
\end{points}
\plain{RMS 和 EDF 的本质区别是“优先级按周期固定下来”还是“优先级随当前截止期限动态变化”。}
\exam{常考比较题：谁是静态优先级、谁是动态优先级、谁通常能支持更高利用率。}


\concept{RMS 可调度上界}
\begin{points}
  \item RMS 采用抢占式静态优先级策略，更短周期的任务优先级更高。
  \item 虽然总利用率 $\sum t_i/p_i \le 1$ 是必要条件，但对 RMS 来说还不总是充分。
  \item 调度 $N$ 个周期任务时，最坏情况下 CPU 利用率上界为 $N(2^{1/N}-1)$。
  \item 因此有些“总利用率不到 100\%”的任务集，按 RMS 仍可能不可调度。
\end{points}
\plain{RMS 不是只看总利用率小于 1 就一定行，它还有更严格的保守上界。}
\exam{常考为什么 RMS 下利用率小于 1 也可能错过截止期限，以及两任务时上界约为 83\%。}


\concept{比例分享调度与 POSIX 实时调度}
\begin{points}
  \item \term{比例分享调度}强调按预先约定的比例在多个任务、用户或任务组之间分配 CPU 份额，而不只是按先来先服务或固定优先级抢占。
  \item 公平分享调度 fair-share scheduling 可看作比例分享思想的一种实现，更关注“每个用户/用户组分到多少 CPU”，而不只是单个进程。
  \item \term{POSIX 实时调度}常见策略包括 \code{SCHED\_FIFO} 和 \code{SCHED\_RR}：前者按固定优先级先来先服务，后者在同优先级任务间再配合时间片轮转。
  \item 这类实时策略通常服务于实时优先级任务，强调高优先级先运行，并允许更高优先级任务抢占当前任务。
\end{points}
\plain{比例分享关心“CPU 份额怎么分”，POSIX 实时调度关心“同一实时优先级下是一直跑到底，还是轮流给时间片”。}
\exam{PPT 明确点了 \code{SCHED\_FIFO} 和 \code{SCHED\_RR}，选择题里常考两者区别：一个同优先级不轮转，一个同优先级按时间片轮转。}


\concept{多处理器调度}
\begin{points}
  \item 非对称多处理调度中，一个主处理器负责调度，其他处理器执行任务。
  \item 对称多处理调度中，各处理器都可参与调度，常见挑战是负载均衡和缓存亲和性。
  \item 处理器亲和性指尽量让进程继续在原 CPU 上运行，以利用缓存中的数据。
\end{points}
\plain{多核下不只要挑进程，还要挑“放到哪颗 CPU 上跑”才划算。}
\exam{常考非对称/对称多处理、亲和性与负载均衡之间的矛盾。}


\concept{亲和性与负载均衡}
\begin{points}
  \item \term{处理器亲和性}强调尽量让进程留在原 CPU 上运行，以减少缓存失效、TLB 失效等迁移代价。
  \item 亲和性可分为 soft affinity 和 hard affinity。
  \item \term{负载均衡}强调让各 CPU 工作负载尽量均匀，常见实现有 push migration 和 pull migration。
  \item 亲和性与负载均衡常有矛盾：留在原 CPU 更省缓存代价，但迁移到空闲 CPU 更能避免忙闲不均。
\end{points}
\plain{这类题别只背名词，核心矛盾就是“少迁移更省缓存”与“适当迁移更均衡负载”之间的取舍。}
\exam{常考 soft/hard affinity、push/pull migration 的含义，以及亲和性和负载均衡为什么可能冲突。}


\concept{多处理器调度算法}
\begin{points}
  \item \term{负载共享}：系统维护全局就绪队列，空闲处理器从中取线程运行，优点是处理器不易闲置，缺点是公共队列可能成为瓶颈。
  \item \term{群调度}：把一组相关线程同时调度到一组处理器上运行，适合同步紧密的并行任务。
  \item \term{专用处理器分配}：给一个应用固定分配一组处理器直到其运行结束，追求高并行度但可能牺牲单个处理器利用率。
  \item \term{动态调度}：OS 负责在应用间分配处理器，应用自己决定其获得的处理器上哪些线程运行，灵活但实现复杂。
\end{points}
\plain{多处理器算法题重点看“用一个全局池分配”，还是“按线程组一起分配”，还是“直接包核给应用”。}
\exam{常考几种多处理器调度算法的基本思想、优缺点和适用场景。}


\concept{线程调度}
\begin{points}
  \item 用户级线程由线程库调度，内核只看到进程。
  \item 内核级线程由内核调度，可以在多个处理器上并行。
  \item 线程调度更强调轻量切换和同一进程内线程之间的同步关系。
\end{points}
\plain{线程调度这节重点不在重新定义线程，而在分清“用户级线程谁在调度、内核级线程谁在调度”，以及这会怎样影响阻塞和并行能力。}
\exam{常考进程/线程区别、用户级线程和内核级线程优缺点、多对一/一对一/多对多模型。}


\concept{典型系统调度实例}
\begin{points}
  \item \term{UNIX System V}：采用动态优先级多级队列调度，CPU 用得越久优先级越低，阻塞等待后返回的交互/I/O 型进程优先级相对更高。
  \item \term{Solaris}：基于优先级的线程调度，不同调度类（如实时、分时、公平分享）可采用不同策略。
  \item \term{Linux}：现代普通任务默认使用 CFS，实时任务常采用符合 POSIX 的 \code{SCHED\_FIFO}/\code{SCHED\_RR}。
  \item \term{Windows}：以线程为调度对象，采用动态优先级、抢占式、多级反馈队列思想，并对交互/I/O 型线程做短暂优先级提升。
\end{points}
\plain{这部分不要求把各系统实现细节全背下来，更重要的是抓住“UNIX 动态优先级、Linux CFS、Windows 线程调度、Solaris 多调度类”这几个标签。}
\exam{若简答题问“现代系统调度器有什么共同趋势”，常可答动态优先级、偏向交互任务、支持实时类和多处理器环境。}


\concept{调度算法的评估方法}
\begin{points}
  \item \term{确定性模型}：给定一组确定的进程到达时间和服务时间，直接计算不同算法下的等待时间、周转时间等指标，适合课堂示例和手工比较。
  \item \term{排队模型}：把进程到达和服务时间看成概率分布，用排队论分析 CPU 利用率、平均等待时间、平均队列长度等期望指标。
  \item \term{仿真法}：用程序模拟调度过程，可用随机数据或真实系统跟踪数据作为输入，结果更接近实际，但实现和统计成本更高。
  \item 三者大致对应“最简单但最理想化”“数学分析更强但依赖分布假设”“最贴近现实但成本最高”的递进关系。
\end{points}
\plain{这一节不是新调度算法，而是在讲“怎么比较算法好坏”：小题目手算常用确定性模型，系统级研究更常见排队分析和仿真。}
\exam{PPT 在章末单独列了评估方法，选择题常考三者各自优缺点：确定性模型快但依赖输入，排队模型依赖分布假设，仿真法更真实但成本高。}
